\documentclass[12pt,a4paper]{article}
\usepackage[UTF8]{ctex}
\usepackage{amsmath,amssymb,amsthm}
\usepackage{geometry}

\geometry{margin=2.5cm}

\title{混合运动-力控制的术语说明}
\author{第11章补充说明}
\date{}

\begin{document}

\maketitle

\section{问题提出}

用户问："所以这东西叫做力位混合控制么？"

答案是：\textbf{基本正确，但更准确的术语是"混合运动-力控制"}。

\section{术语对比}

\subsection{英文术语}

\textbf{标准术语}：
\begin{itemize}
\item \textbf{Hybrid Motion-Force Control}
\item 直译：混合运动-力控制
\end{itemize}

\textbf{其他常见说法}：
\begin{itemize}
\item Hybrid Position-Force Control（混合位置-力控制）
\item Hybrid Control（混合控制，简称）
\end{itemize}

\subsection{中文术语}

\textbf{标准术语}：
\begin{itemize}
\item \textbf{混合运动-力控制}
\item 这是最准确的翻译
\end{itemize}

\textbf{其他常见说法}：
\begin{itemize}
\item \textbf{力位混合控制}（常用，但不够准确）
\item 混合位置-力控制
\item 混合控制（简称）
\end{itemize}

\section{为什么叫"混合"控制？}

\subsection{基本含义}

\textbf{"混合"的含义}：
\begin{itemize}
\item 同时使用两种控制方法
\item 在不同方向上分别使用运动控制和力控制
\item 将两种控制策略"混合"在一起
\end{itemize}

\subsection{具体说明}

\textbf{混合运动-力控制器}：
\begin{equation}
\tau = J_b^T(\theta)\Bigg[P(\theta)\text{（运动控制）} + (I - P(\theta))\text{（力控制）} + \text{补偿项}\Bigg]
\end{equation}

\textbf{组成}：
\begin{enumerate}
\item \textbf{运动控制部分}：在自由方向上控制运动
\item \textbf{力控制部分}：在约束方向上控制力
\item \textbf{混合方式}：通过投影矩阵$P$和$I-P$将两者分离并组合
\end{enumerate}

\section{术语的准确性}

\subsection{"力位混合控制" vs "混合运动-力控制"}

\textbf{"力位混合控制"}：
\begin{itemize}
\item \textbf{优点}：简洁，常用
\item \textbf{缺点}：
\begin{itemize}
\item "位"（位置）不够准确，因为不仅仅是位置，还包括速度、加速度等
\item 应该用"运动"（motion）更准确
\end{itemize}
\end{itemize}

\textbf{"混合运动-力控制"}：
\begin{itemize}
\item \textbf{优点}：
\begin{itemize}
\item 更准确，因为"运动"包括位置、速度、加速度等
\item 与英文术语一致
\item 更专业
\end{itemize}
\item \textbf{缺点}：稍长一些
\end{itemize}

\subsection{为什么"运动"比"位置"更准确？}

\textbf{运动控制包括}：
\begin{itemize}
\item 位置控制：$x_d$
\item 速度控制：$\dot{x}_d$
\item 加速度控制：$\ddot{x}_d$
\item 轨迹跟踪：$x_d(t)$
\end{itemize}

\textbf{位置控制只是运动控制的一种}：
\begin{itemize}
\item 位置控制：只控制位置
\item 运动控制：包括位置、速度、加速度等
\end{itemize}

\textbf{在混合控制中}：
\begin{itemize}
\item 我们控制的是"运动"（motion），不仅仅是"位置"（position）
\item 例如：控制速度$v_x$，而不仅仅是位置$x$
\end{itemize}

\section{术语使用建议}

\subsection{正式场合}

\textbf{推荐使用}：
\begin{itemize}
\item \textbf{混合运动-力控制}（Hybrid Motion-Force Control）
\item 最准确，最专业
\end{itemize}

\subsection{非正式场合}

\textbf{可以使用}：
\begin{itemize}
\item \textbf{力位混合控制}
\item 简洁，易于理解
\item 虽然不够准确，但广泛使用
\end{itemize}

\subsection{其他说法}

\textbf{也可以使用}：
\begin{itemize}
\item 混合位置-力控制（Hybrid Position-Force Control）
\item 混合控制（Hybrid Control，简称）
\end{itemize}

\section{相关术语}

\subsection{控制方法分类}

\begin{enumerate}
\item \textbf{运动控制}（Motion Control）：
\begin{itemize}
\item 只控制运动，不控制力
\item 适用于自由空间
\end{itemize}

\item \textbf{力控制}（Force Control）：
\begin{itemize}
\item 只控制力，不控制运动
\item 适用于完全约束的环境
\end{itemize}

\item \textbf{混合运动-力控制}（Hybrid Motion-Force Control）：
\begin{itemize}
\item 在不同方向上分别控制运动和力
\item 适用于部分约束的环境
\end{itemize}

\item \textbf{阻抗控制}（Impedance Control）：
\begin{itemize}
\item 控制机器人的阻抗特性（质量、弹簧、阻尼）
\item 感测运动，产生力
\end{itemize}

\item \textbf{导纳控制}（Admittance Control）：
\begin{itemize}
\item 控制机器人的导纳特性
\item 感测力，产生运动
\end{itemize}
\end{enumerate}

\subsection{术语关系}

\textbf{混合运动-力控制}是：
\begin{itemize}
\item 运动控制和力控制的\textbf{组合}
\item 在不同方向上分别使用两种控制方法
\item 通过投影矩阵实现"混合"
\end{itemize}

\section{总结}

\subsection{术语对比表}

\begin{center}
\begin{tabular}{|c|c|c|}
\hline
\textbf{中文术语} & \textbf{英文术语} & \textbf{准确性} \\
\hline
混合运动-力控制 & Hybrid Motion-Force Control & \checkmark 最准确 \\
\hline
力位混合控制 & Hybrid Position-Force Control & $\sim$ 常用但不完全准确 \\
\hline
混合位置-力控制 & Hybrid Position-Force Control & $\sim$ 可接受 \\
\hline
混合控制 & Hybrid Control & $\sim$ 简称 \\
\hline
\end{tabular}
\end{center}

\subsection{推荐使用}

\begin{itemize}
\item \textbf{正式文档}：使用"混合运动-力控制"
\item \textbf{日常交流}：可以使用"力位混合控制"（简洁易懂）
\item \textbf{学术论文}：使用"混合运动-力控制"或"Hybrid Motion-Force Control"
\end{itemize}

\subsection{关键理解}

无论使用哪个术语，关键是要理解：
\begin{enumerate}
\item \textbf{混合}：同时使用两种控制方法
\item \textbf{在不同方向上}：不是在同一个方向上
\item \textbf{分别控制}：运动控制和力控制分别作用于不同的方向
\item \textbf{通过投影}：使用投影矩阵实现分离和组合
\end{enumerate}

\end{document}

